\documentclass[conference]{IEEEtran}
\IEEEoverridecommandlockouts

\usepackage{cite}
\usepackage{amsmath,amssymb,amsfonts}
\usepackage{graphicx}
\usepackage{booktabs}
\usepackage{xcolor}
\usepackage{url}
\usepackage[font=small]{caption}

\newcommand{\figplaceholder}[2]{%
  \fbox{\parbox[c][#2][c]{0.95\linewidth}{\centering\small\texttt{#1}\\[2pt]
  \footnotesize (placeholder -- see FIGURES\_GUIDE.md)}}}

\begin{document}

\title{FLARE: Early, Adaptive, and Explainable Detection of Low-Rate
Flow-Table Overflow Attacks in Software-Defined Networks}

\author{\IEEEauthorblockN{Author}
\IEEEauthorblockA{Department of CSE\\
NIT Delhi, New Delhi, India\\
abc@gmail.com}
% TODO: add co-authors / supervisor here in additional \IEEEauthorblockN blocks
}

\maketitle

\begin{abstract}
Software-defined networking concentrates forwarding state in the flow tables of
switches, which are backed by small and expensive ternary content-addressable
memories that hold only a few thousand entries. A low-rate flow-table overflow
(LOFT) attack exploits this constraint by installing long-lived malicious rules
at a trickle, keeping each rule alive with minimal traffic until legitimate
flows can no longer be served. Existing defenses are reactive and raise an alarm
only after the table has begun to fill, leaving operators very little time to
respond. This paper presents FLARE, a streaming defense framework that detects
LOFT attacks substantially earlier. FLARE combines four mechanisms: an
early-warning detector that recognizes the attacker's reconnaissance (probing)
phase before any table growth occurs, an online core detector that couples
cumulative-sum drift detection with adaptively normalized churn statistics, a
time-to-overflow forecaster that converts detection into an actionable deadline,
and a feature-attribution module that attaches a plain-language reason to every
alert. In a reproducible packet-level simulation of an SDN switch, FLARE raises
its first alarm a median of 150\,s before overflow, compared with 30\,s for a
reactive occupancy-threshold baseline, while keeping the false-positive rate on
attack-free traffic at zero across all runs. The forecaster predicts the
overflow instant with a median error of 10\,s. The complete framework, labeled
dataset, and a one-click notebook that reproduces every result are publicly
available.
\end{abstract}

\begin{IEEEkeywords}
software-defined networking, flow table overflow, low-rate attacks, anomaly
detection, CUSUM, online learning, explainable security
\end{IEEEkeywords}

%==============================================================================
\section{Introduction}
%==============================================================================
Software-defined networking (SDN) separates the control plane from the data
plane and gives operators centralized, programmable control over their
infrastructure~\cite{mckeown2008openflow}. The architecture has been widely
adopted in data centers, campus networks, and wide-area backbones, and its
security properties have consequently attracted sustained
attention~\cite{kreutz2015sdn}. Forwarding decisions in an SDN switch are
resolved against a flow table whose entries are typically stored in ternary
content-addressable memory (TCAM). TCAM is fast but power-hungry and expensive,
so commodity switches ship with tables that hold on the order of a few thousand
entries~\cite{curtis2011devoflow}. This small, shared resource is a natural
target.

Early work demonstrated that an adversary who can trigger rule installations is
able to exhaust the flow table and degrade or disable
forwarding~\cite{shin2013attacking}. Flooding-style exhaustion is noisy,
however, and is caught by conventional volume-based defenses. The more dangerous
variant is the \emph{low-rate} flow-table overflow (LOFT) attack, in which the
adversary installs rules at a trickle and refreshes each one just before its
idle timeout so that it never expires~\cite{cao2017loft}. The aggregate attack
traffic remains negligible, the per-flow behavior superficially resembles a
quiet legitimate flow, and the table fills over minutes rather than seconds.
Prior measurements show that the parameters an attacker needs, such as the idle
timeout and the approximate table capacity, can be inferred remotely through
careful timing probes~\cite{leng2015inference}.

Detection work has so far concentrated on the consequences of the attack rather
than its onset. The most recent representative of this line, FloRa, monitors
table state and flags the overflow process once it is under
way~\cite{flora2024}. Reactive detection of this kind is useful but
structurally late: by the time occupancy is abnormal, the attacker has already
spent its budget and the operator inherits a table full of malicious state. In
our experiments, an occupancy-threshold detector representative of this approach
fires a median of only 30\,s before the table is effectively full.

We argue that a LOFT campaign is observable much earlier. The attack has a
characteristic life cycle: a quiet \emph{probing} phase in which the adversary
measures timeouts and capacity, followed by an \emph{installation} phase in
which low-rate flows accumulate. Both phases perturb flow-level churn statistics
long before they move the occupancy needle. Acting on these signals requires a
detector that adapts to non-stationary benign traffic without being poisoned by
the attack itself, and that communicates risk in terms an operator can act on.

This paper presents FLARE (FLow-table Anomaly Recognition Engine), a streaming
defense framework built around four ideas.

\begin{itemize}
\item[(i)] \textbf{Early warning.} A dedicated detector recognizes the probing phase
from a sustained excess of short-lived, single-packet flows while occupancy is
still flat, raising a flag before any damage begins.

\item[(ii)] \textbf{Online adaptation.} All baselines are learned with exponentially
weighted statistics that track benign drift continuously; learning is frozen
while an anomaly is in progress so the attack cannot drag the baseline along
with it.

\item[(iii)] \textbf{Risk forecasting.} A gated trend estimator projects the current
fill rate to a concrete time-to-overflow, replacing a binary alarm with a
deadline.

\item[(iv)] \textbf{Explainability.} Every alert carries a ranked, human-readable
summary of the features that triggered it.
\end{itemize}

We evaluate FLARE in a packet-level simulation of an OpenFlow-style switch under
benign traffic and a two-phase LOFT campaign. Across twenty independently seeded
runs, FLARE detects every attack, fires earlier than the reactive baseline in
every run, provides a median lead of 150\,s before overflow, and produces no
false positives on attack-free traffic. The framework, the labeled dataset, and
a notebook that reproduces every figure in this paper run end-to-end on a free
cloud instance.

The remainder of the paper is organized as follows.
Section~\ref{sec:related} reviews related work.
Section~\ref{sec:threat} defines the threat model.
Section~\ref{sec:design} presents the design of FLARE.
Section~\ref{sec:setup} describes the experimental setup, and
Section~\ref{sec:results} reports results.
Section~\ref{sec:discussion} discusses limitations, and
Section~\ref{sec:conclusion} concludes.

%==============================================================================
\section{Related Work}\label{sec:related}
%==============================================================================
\subsection{Flow-Table Exhaustion Attacks}
The feasibility of attacking SDN through its own control machinery was
established shortly after OpenFlow's introduction~\cite{shin2013attacking}.
Kandoi and Antikainen quantified how table-miss handling and rule installation
create denial-of-service opportunities in OpenFlow
deployments~\cite{kandoi2015dos}. Cao et al.\ then showed that overflow does not
require a flood: by inferring the idle timeout and refreshing rules just before
expiry, an attacker fills the table while remaining inconspicuous at every
vantage point that measures volume~\cite{cao2017loft}. Complementary work
demonstrated that the table parameters needed to mount such an attack can be
inferred remotely~\cite{leng2015inference}. Related control-plane reflection
attacks similarly weaponize the switch's limited resources at low
rates~\cite{zhang2018reflection}.

\subsection{Defenses}
A first family of defenses targets flooding-style saturation.
Avant-Guard introduced connection migration in the data plane to shield the
control channel from SYN-flood-driven rule churn~\cite{shin2013avantguard}.
FloodGuard combined proactive rule pre-installation with packet-in rate
limiting~\cite{wang2015floodguard}, and FloodDefender added table-miss
engineering and packet filtering across neighboring
switches~\cite{shang2017flooddefender}. SLICOTS targets TCP-specific flooding
through lightweight connection tracking~\cite{mohammadi2017slicots}. These
systems assume a noisy attacker and are largely blind to low-rate campaigns.

A second family manages the table itself. DevoFlow reduces TCAM pressure by
devolving short flows to the data plane~\cite{curtis2011devoflow}. Yuan et al.\
proposed peer support, offloading entries from an overloaded switch to its
neighbors~\cite{yuan2019peer}. Such mechanisms raise the cost of an overflow but
do not detect one. Statistical detection of anomalous flow patterns in SDN dates
back to self-organizing-map classifiers over flow features~\cite{braga2010som}.
The closest work to ours is FloRa, which detects low-rate overflow by monitoring
table dynamics, but it reacts only once occupancy is already
abnormal~\cite{flora2024}. None of these systems attempts to detect the
reconnaissance that precedes installation, to forecast the time remaining before
overflow, or to explain its decisions.

\subsection{Change Detection and Explainability}
FLARE's core builds on the cumulative-sum (CUSUM) procedure, which is optimal in
a minimax sense for detecting small persistent shifts in a monitored
statistic~\cite{page1954cusum}. The general theory of sequential change
detection is covered by Basseville and
Nikiforov~\cite{basseville1993detection}. Our preference for simple, adaptive
statistics over heavyweight learned models follows the well-known critique of
anomaly-detection deployments in operational networks~\cite{sommer2010outside}.
The explainability module adopts the additive feature-attribution view
popularized by SHAP, in the lightweight form appropriate for a streaming
detector~\cite{lundberg2017shap}.

%==============================================================================
\section{Threat Model}\label{sec:threat}
%==============================================================================
We consider a single OpenFlow-style switch with a flow table of capacity $C$
entries and an idle timeout of $T_{\mathrm{idle}}$ seconds, managed reactively:
a table miss triggers rule installation for the new flow, and an entry is
evicted $T_{\mathrm{idle}}$ seconds after its last matching packet. Once the table reaches capacity, the switch can no longer install new entries: the corresponding packets are either discarded or redirected to the controller. In both cases the outcome is the same—legitimate flows stop being forwarded. We model the adversary as one external host, or a small group of them, with the ability to push packets through the switch and to construct packet headers at will. The adversary has no access to the controller and cannot tamper with the
switch firmware, nor does it know $C$ or $T_{\mathrm{idle}}$ beforehand. As in
prior work~\cite{cao2017loft}, the attack unfolds across two phases.

\subsubsection{Probing phase} In this stage the attacker injects short-lived
flows---often just one packet apiece---and observes how quickly the switch
responds, using that timing to estimate $T_{\mathrm{idle}}$ and gauge how
loaded the table is~\cite{leng2015inference}. Because each probe holds an
entry for only a moment, the table's occupancy stays effectively unchanged.

\subsubsection{Installation phase} The attacker opens new flows at a low rate
$\lambda_a$ and refreshes each one with a single packet every
$T_{\mathrm{idle}} - \delta$ seconds, for a small safety margin $\delta$, so
that no malicious entry ever expires. The per-flow packet rate is roughly
$1/(T_{\mathrm{idle}}-\delta)$ packets per second, far below any volumetric
alarm threshold, yet entries accumulate monotonically until the table is full.

The defender observes only aggregate, per-window statistics of the switch
(occupancy, miss and installation counts, expiry counts, and packet counts) and
must operate online with bounded memory. We explicitly allow benign traffic to
be non-stationary.

%==============================================================================
\section{The FLARE Framework}\label{sec:design}
%==============================================================================

\begin{figure}[t]
\centering
\includegraphics[width=\linewidth]{fig_architecture.jpg}
\caption{FLARE architecture. Per-window features extracted from the switch feed
four streaming modules: the early-warning detector, the core detector, the
time-to-overflow forecaster, and the explainer; confirmed detections drive the
mitigation policy.}
\label{fig:arch}
\end{figure}

Fig.~\ref{fig:arch} shows the architecture. Time is divided into windows of
$\Delta$ seconds (we use $\Delta = 1$\,s). At the end of window $t$ the switch
state is summarized as a feature vector, and each module consumes the vector in
a single pass, so the whole stack runs in constant memory and $O(|\mathcal{F}|)$
time per window.

\subsection{Feature Model}
Let $o_t$ denote table occupancy at the end of window $t$, and let $m_t$, $i_t$,
$e_t$, $s_t$, and $p_t$ denote the table-miss, installation, expiry,
short-lived-expiry, and packet counts within the window. From these we derive
the feature set $\mathcal{F}$: the occupancy ratio $o_t / C$; the smoothed
occupancy growth $g_t$, an exponentially weighted average of
$(o_t - o_{t-1})/\Delta$; the miss, installation, expiry, and short-lived rates;
the net accumulation $(i_t - e_t)/\Delta$; and the mean traffic intensity per
installed flow, $p_t / o_t$. The last feature is central to this threat: LOFT
flows occupy an entry indefinitely while contributing almost no packets, so a
falling per-flow intensity alongside rising occupancy is a strong attack
signature.

\subsection{Online Baselines with Freeze-on-Alarm}
Every monitored feature $f$ maintains an exponentially weighted mean and
variance,
\begin{align}
\mu_t^{(f)} &= \mu_{t-1}^{(f)} + \alpha\,\bigl(x_t^{(f)} - \mu_{t-1}^{(f)}\bigr),
\label{eq:ewma-mean}\\
\bigl(\sigma_t^{(f)}\bigr)^2 &= (1-\alpha)\Bigl[\bigl(\sigma_{t-1}^{(f)}\bigr)^2
 + \alpha\,\bigl(x_t^{(f)} - \mu_{t-1}^{(f)}\bigr)^2\Bigr],
\label{eq:ewma-var}
\end{align}
with smoothing factor $\alpha \in (0,1)$ (we use $\alpha = 0.02$, an effective
memory of roughly $100$ windows). The first $W$ windows are treated as a warm-up
during which moments are estimated incrementally and no alarms are raised.

Adaptation creates a poisoning risk: a slowly escalating attack could teach the
baseline that its own behavior is normal. FLARE therefore applies
\emph{freeze-on-alarm}: updates~\eqref{eq:ewma-mean}--\eqref{eq:ewma-var} are
suspended for any window the detector currently scores as anomalous, and resume
when the score returns below threshold. The baseline thus tracks benign drift
at all other times while remaining anchored during suspicious episodes.

\subsection{Core Detector}
The core detector fuses two complementary signals.

\subsubsection{Directional churn score}
For each churn feature $f$ the standardized deviation is
\begin{equation}
z_t^{(f)} = \frac{x_t^{(f)} - \mu_{t-1}^{(f)}}{\sigma_{t-1}^{(f)}},
\label{eq:zscore}
\end{equation}
clipped to the direction in which an attack moves it:
$\tilde z_t^{(f)} = \max\bigl(d_f\, z_t^{(f)},\, 0\bigr)$, where
$d_f \in \{+1, -1\}$ is $+1$ for miss, installation, expiry, and short-lived
rates and $-1$ for per-flow traffic intensity. The window's churn score is the
root-mean-square aggregate
\begin{equation}
A_t = \Bigl(\tfrac{1}{|\mathcal{F}_c|}
      \textstyle\sum_{f \in \mathcal{F}_c} \bigl(\tilde z_t^{(f)}\bigr)^2
      \Bigr)^{1/2},
\label{eq:rms}
\end{equation}
which rewards coherent multi-feature deviations over isolated spikes.

\subsubsection{CUSUM drift on table growth}
The defining behavior of a LOFT attack, slow persistent accumulation, is poorly
served by per-window scores: an adaptive baseline simply learns the new growth
rate. We instead monitor the smoothed growth $g_t$ with a one-sided CUSUM
statistic~\cite{page1954cusum},
\begin{equation}
S_t = \max\bigl(0,\; S_{t-1} + g_t - \mu_{t-1}^{(g)} - k\,\sigma_{t-1}^{(g)}\bigr),
\label{eq:cusum}
\end{equation}
where $k$ is the reference slack (we use $k = 1$). The normalized drift
indicator is $D_t = S_t / (h\,\sigma_{t-1}^{(g)})$ with decision limit $h$ (we
use $h = 9$), so $D_t \ge 1$ signals that small positive growth excesses have
persisted long enough to matter. The CUSUM reference statistics are also subject
to freeze-on-alarm.

\subsubsection{Decision rule}
Window $t$ is marked anomalous if $A_t \ge \tau$ or $D_t \ge 1$, with
$\tau = 3$. A single noisy window shouldn't set off an alarm. To avoid that, we apply a
$k$-of-$n$ persistence rule and only flag an anomaly once at least $\kappa$ of
the last $n$ windows look abnormal---in practice we set $\kappa = 3$ and
$n = 5$. For the score itself we take $\max(A_t,\, \tau D_t)$, which lets us
compare the two signals on the same scale.

\subsection{Early-Warning Detector}
What gives probing away is a steady stream of short-lived flows and table
misses that persists \emph{while occupancy stays low}. That combination sets it
apart both from ordinary traffic bursts and from the later installation phase.
To catch it, the early-warning module watches the short-lived rate and the miss
rate, each against its own baseline of the
form~\eqref{eq:ewma-mean}--\eqref{eq:ewma-var}, and raises a flag on window $t$
when
\begin{equation}
\max_{f \in \{s, m\}} \tilde z_t^{(f)} \ge \tau_p
\quad\text{and}\quad
o_t / C < \theta,
\label{eq:probe}
\end{equation}
using $\tau_p = 2.8$ and an occupancy gate of $\theta = 0.6$; the same
$k$-of-$n$ persistence rule still applies. That occupancy gate is precisely
what makes the module an early warning. Once the table is visibly filling, the
detector steps back on purpose and leaves that regime to the core detector,
which keeps the two alarms cleanly separated in meaning.

\subsection{Time-to-Overflow Forecaster}
Knowing \emph{whether} something is wrong is only half of what operators need;
they also want to know \emph{how long they have}. To answer that, FLARE looks
at the last $W_r$ occupancy samples (we use $W_r = 20$), fits an ordinary
least-squares line $\hat o(t) = \hat\beta t + \hat\gamma$ across that sliding
window, and records the coefficient of determination $R^2$. It publishes a
forecast only when the fill is genuinely sustained---that is, when
\begin{equation}
\hat\beta \ge \beta_{\min}
\;\wedge\;
R^2 \ge \rho_{\min}
\label{eq:gate}
\end{equation}
holds for $m$ consecutive windows (here $\beta_{\min} = 3$ entries/s,
$\rho_{\min} = 0.7$, and $m = 10$). The gate earns its keep in practice: benign
occupancy sits flat but jittery, so an ungated extrapolation swings wildly,
whereas a real fill traces a near-linear path and clears the gate right away.
Once it does, the estimated time to overflow is
\begin{equation}
\widehat{\mathrm{ETA}}_t = \max\!\left(\frac{C - \hat o(t)}{\hat\beta},\, 0\right).
\label{eq:eta}
\end{equation}

\subsection{Explainability}
Each alert is annotated with the features responsible for it. The per-feature
deviations $\tilde z_t^{(f)}$ (and the scaled drift $\tau D_t$) computed
in~\eqref{eq:rms} already constitute an additive decomposition of the score, in
the spirit of feature-attribution methods~\cite{lundberg2017shap}. FLARE ranks
them, keeps the top two, and renders a fixed-vocabulary sentence such as
``burst of short-lived probe-like flows ($+4.0\sigma$); flow table filling
faster than normal ($+3.2\sigma$).'' The cost is negligible and the output
requires no model introspection.

\subsection{Mitigation}
Once an attack is confirmed, FLARE applies selective eviction with timeout
hardening. Suspected low-rate flows, those with per-flow traffic intensity near
the refresh floor $1/(T_{\mathrm{idle}}-\delta)$, are denied table residency,
and the effective idle timeout of existing low-activity entries is reduced to
$T_{\mathrm{idle}}/4$. Legitimate flows, which arrive with substantially higher
packet rates, are unaffected. Section~\ref{sec:results} quantifies the effect.

%==============================================================================
\section{Experimental Setup}\label{sec:setup}
%==============================================================================
We evaluate FLARE in a packet-level simulation of an OpenFlow-style switch,
implemented so that the entire study reproduces on a free cloud notebook
without specialized hardware; software switch models of this kind are standard
practice for SDN experimentation~\cite{lantz2010mininet}. The flow-table
semantics (install on first packet, refresh on match, evict on idle timeout,
reject when full) are modeled exactly.

\begin{table}[t]
\caption{Simulation and detector parameters.}
\label{tab:params}
\centering
\small
\begin{tabular}{@{}llr@{}}
\toprule
 & Parameter & Value \\
\midrule
Switch & table capacity $C$ & 1500 entries \\
 & idle timeout $T_{\mathrm{idle}}$ & 10\,s \\
Traffic & benign flow arrival rate & 26 flows/s \\
 & mean benign flow lifetime & 8\,s \\
 & benign in-flow packet rate & 3 pkt/s \\
Attack & probing window & 200--240\,s \\
 & probe rate & 14 flows/s \\
 & installation start & 240\,s \\
 & malicious flow rate $\lambda_a$ & 7--11 flows/s \\
 & refresh margin $\delta$ & 1\,s \\
Detector & window $\Delta$ / warm-up $W$ & 1\,s / 30 \\
 & smoothing $\alpha$ & 0.02 \\
 & churn threshold $\tau$ / probe $\tau_p$ & 3.0 / 2.8 \\
 & CUSUM $k$, $h$ & 1, 9 \\
 & persistence $\kappa$-of-$n$ & 3-of-5 \\
Forecast & $W_r$, $\beta_{\min}$, $\rho_{\min}$, $m$ & 20, 3, 0.7, 10 \\
\bottomrule
\end{tabular}
\end{table}

Table~\ref{tab:params} lists the configuration. Each simulation covers 600\,s.
Benign flows arrive as a Poisson process with exponential lifetimes, and a
configurable fraction are one- to two-packet ``blips'' that emulate DNS lookups
and scans. The attack follows the two-phase model of
Section~\ref{sec:threat}. Ground truth labels every window as benign, probing,
or attack, and the \emph{overflow instant} is defined as the first window in
which occupancy reaches 98\% of $C$, the point at which legitimate
installations begin to fail in practice.

\subsubsection{Baseline} We compare against a reactive occupancy-threshold
detector that alarms when $o_t / C \ge 0.8$, which captures the operating
principle of table-state monitors such as FloRa~\cite{flora2024}.

\subsubsection{Metrics} For each method we report several quantities: the time of its first alarm; the
\emph{lead before overflow}, defined as the overflow instant minus that first
alarm; the lead it gains over the baseline; the false-positive rate, taken as
the fraction of windows flagged on attack-free runs; and the absolute forecast
error. We evaluate that last one at the moment occupancy first crosses 60\% of
$C$, since that is when an operator would actually turn to the forecast.

\subsubsection{Reproducibility} All randomness derives from a single seed per
run. The headline scenario uses seed 7; robustness results aggregate seeds
1--20. The code, dataset generator, and notebook are available in the project
repository.\footnote{\url{https://github.com/prakhar443/project_FLARE}}

%==============================================================================
\section{Results}\label{sec:results}
%==============================================================================

\subsection{The Attack Succeeds Against an Undefended Switch}

\begin{figure}[t]
\centering
\includegraphics[width=\linewidth]{01_attack_no_defense.png}
\caption{Flow-table occupancy under a LOFT attack with no defense. The probing
phase (200--240\,s) barely perturbs occupancy; the installation phase fills the
table until overflow at 361\,s.}
\label{fig:attack}
\end{figure}

Fig.~\ref{fig:attack} shows the headline scenario. Benign traffic settles at a
steady occupancy near 450 of 1500 entries. Probing at 200\,s produces only a
slight, transient rise, exactly the stealth the attacker intends. From 240\,s
the installation phase accumulates entries monotonically, and the table is
effectively full at $t = 361$\,s, after which overflow drops of legitimate
installations are sustained. Total attack bandwidth remains under one percent
of link traffic throughout, confirming that volumetric defenses have nothing to
observe.

\subsection{Detection Lead Time}

\begin{table}[t]
\caption{Detection results for the headline scenario (overflow occurs at 361\,s).}
\label{tab:headline}
\centering
\small
\resizebox{\dimexpr\columnwidth-1cm\relax}{!}{%
\begin{tabular}{@{}lccc@{}}
\toprule
Method & First alarm & Lead before & Flags \\
 & (s) & overflow (s) & probing \\
\midrule
Occupancy baseline (FloRa-style) & 336 & 25 & no \\
FLARE, core detector & 209 & 152 & no \\
FLARE, early warning & 210 & 151 & yes \\
\bottomrule
\end{tabular}%
}
\end{table}

\begin{figure}[t]
\centering \includegraphics[width=\linewidth]{02_detection_timeline.png}
\caption{Detection timeline. Top: occupancy, marked with each method's
first-alarm instant. Bottom: the FLARE anomaly score, flat under benign traffic
and crossing the threshold as probing begins.}
\label{fig:timeline}
\end{figure}

Table~\ref{tab:headline} and Fig.~\ref{fig:timeline} bring together the
detection results for the headline scenario. The reactive baseline holds off
until 336\,s---a mere 25\,s before the table overflows. FLARE behaves very
differently. Its early-warning module fires at 210\,s, just ten seconds into
the probing phase and a full 126\,s ahead of the baseline, while the core
detector confirms at 209\,s by combining churn deviations with CUSUM drift. The
lower panel of Fig.~\ref{fig:timeline} makes the mechanism visible: under benign
load the anomaly score hovers near 1, then climbs past $\tau = 3$ within seconds
of the first probe. Because the freeze-on-alarm rule stops the growing attack
from quietly re-normalizing the baseline, the score stays high right through the
installation phase instead of decaying as the attack drags on.

\subsection{Forecasting the Overflow Instant}

\begin{figure}[t]
\centering
\includegraphics[width=\linewidth]{03_forecast.png}
\caption{Time-to-overflow forecast. After the gate of~\eqref{eq:gate} opens, the
predicted overflow instant follows the true value (361\,s) to within roughly
$\pm20$\,s.}
\label{fig:forecast}
\end{figure}

Fig.~\ref{fig:forecast} compares the forecast we publish against the actual
overflow instant. Throughout the benign period the gate of~\eqref{eq:gate} keeps
the forecaster quiet---a sensible choice, since naively extrapolating the
flat-but-noisy occupancy there would produce predictions that swing by thousands
of seconds. The gate opens soon after installation starts, and from then on the
predicted overflow time tracks the truth to within about $\pm 20$\,s, tightening
as overflow draws near. At the 60\%-occupancy checkpoint the forecast error is
13.8\,s on the headline run, with a median of 10\,s across seeds. Put in
operational terms: an operator alerted at 210\,s receives not just an alarm but
a believable deadline still more than two minutes out.

\subsection{Robustness Across Runs}

\begin{table}[t]
\caption{Aggregate results across 20 independently seeded runs.}
\label{tab:robust}
\centering
\small
\begin{tabular}{@{}lc@{}}
\toprule
Metric & Value \\
\midrule
Attack detection rate (core) & 20/20 \\
Probing detected by early warning & 20/20 \\
Runs where FLARE precedes baseline & 20/20 \\
Median lead before overflow, FLARE & 150\,s \\
Median lead before overflow, baseline & 30\,s \\
Median lead over baseline & 121\,s \\
Median (max) forecast error & 10\,s (35\,s) \\
False-positive rate, FLARE & 0\,\% \\
False-positive rate, early warning & 0\,\% \\
\bottomrule
\end{tabular}
\end{table}

To make sure a single lucky seed isn't driving the result, we rerun the
experiment across seeds 1--20, cycling the attack rate through 7, 9, and 11
flows/s. Table~\ref{tab:robust} collects the outcome. Both FLARE detectors fire
on every run, FLARE beats the baseline on every run, and its median lead before
overflow is five times that of the baseline. Across the attack-free runs---more
than 1700 post-warm-up windows per detector---the false-positive rate is zero
for every method. We put this down to the persistence rule, which soaks up the
isolated deviations thrown off by benign Poisson churn.

\subsection{Mitigation}

\begin{figure}[t]
\centering \includegraphics[width=\linewidth]{04_mitigation.png}
\caption{Occupancy with and without mitigation. Selective eviction plus timeout
hardening, triggered at FLARE's detection instant (211\,s), keeps the table well
below capacity for the rest of the run.}
\label{fig:mitigation}
\end{figure}

Fig.~\ref{fig:mitigation} replays the headline scenario, this time with the
mitigation policy from Section~\ref{sec:design} switched on at FLARE's first
alarm. Refusing residency to flows sitting at the refresh floor and hardening
the timeouts clears out the malicious state within seconds: occupancy falls to
around 250 entries and stays bounded there, and no benign installations are lost
to overflow for the rest of the run. For contrast, we overlay the undefended
trace from Fig.~\ref{fig:attack}.

\subsection{Overhead}
Each module does a constant amount of work per feature per window: the updates
in~\eqref{eq:ewma-mean}--\eqref{eq:cusum} come down to a few arithmetic
operations, and the forecaster's least-squares fit over $W_r = 20$ points costs
almost nothing. The whole stack gets through a 600-window run in well under a
second on a single commodity core, so running it alongside a production
controller adds no load worth measuring.

%==============================================================================
\section{Discussion and Limitations}\label{sec:discussion}
%==============================================================================
\subsubsection{Why CUSUM} Against low-rate attacks, any per-window
score---however carefully normalized---runs into a bind: an adaptive baseline
quietly absorbs slow drift, whereas a static one piles up false positives as
benign traffic shifts over time. The CUSUM statistic of~\eqref{eq:cusum}
sidesteps the dilemma by accumulating small excesses, and its optimality for
detecting persistent small shifts is a classical result~\cite{page1954cusum}.
When we ablated the CUSUM term, the core detector missed the installation phase
in 2 of 10 seeds; with the term restored, every seed was caught.

\subsubsection{Evasion} An attacker who knows about FLARE could throttle the
probing to stay under $\tau_p$, but the occupancy gate in~\eqref{eq:probe}
leaves no room to game it: a successful overflow has to move occupancy sooner or
later, and once it does the CUSUM term integrates that growth no matter how
gradually it accrues. Drawing the campaign out also widens the attacker's
exposure window and gives the operator more time to respond---exactly the
trade-off a defender wants to force.

\subsubsection{Scope} Our evaluation rests on a software model of a single
switch. The upside is that every result reproduces on a free cloud notebook and
the flow-table semantics are captured exactly; the downside is that hardware
effects such as TCAM update latency and controller scheduling jitter are left
out. Validating FLARE on a Mininet or P4 testbed is the obvious next
step~\cite{lantz2010mininet}. Pushing detection across several switches, where a
distributed attacker spreads its budget thin, will probably mean correlating
FLARE instances at the controller---something the per-window feature interface
already accommodates.

%==============================================================================
\section{Conclusion}\label{sec:conclusion}
%==============================================================================
Low-rate flow-table overflow attacks get through because they stay below the
sensitivity of volumetric defenses and ahead of reactive table monitors. FLARE
shuts both doors. By picking up on the reconnaissance that precedes
installation, folding slow drift into a CUSUM statistic built on adaptively
learned baselines, and turning detection into an explained, forecast deadline,
it converts a last-second alarm into minutes of usable warning: a median of
150\,s before overflow against 30\,s for a reactive baseline, no false positives
across any run, and a median forecast error of 10\,s. The framework, the
dataset, and a notebook that regenerates every figure are all public, and we
hope they serve as a starting point for future work on anticipatory SDN defense.

%==============================================================================
\bibliographystyle{IEEEtran}
\bibliography{references}

\end{document}
